\section{Core Component Implementation}

This section details implementation of core processing components, focusing on key design decisions and code organisation.

\subsection{Frontend: Data Acquisition Module}

The \texttt{DataSource} base class defines a \texttt{stream(chunk\_size)} generator yielding NumPy arrays shaped \texttt{(n\_ants, chunk\_size)}.

The \texttt{SimulatedStream} class generates synthetic signals:

\begin{verbatim}
for each chunk:
    signal = carrier * exp(-2πj * phase_offset) + noise
    yield signal
\end{verbatim}

NumPy broadcasting extends the carrier signal across antennas automatically. Phase offsets are computed once at initialisation.

The \texttt{BatchFileSource} loads data once and yields sequential array views (not copies) to minimise memory overhead.

\subsection{F-Engine: Channelization Implementation}

The F-Engine generates windows at initialisation to amortise cost. Processing iterates over antennas and spectral windows:

\begin{verbatim}
for ant in range(n_ants):
    for spec in range(n_spectra):
        windowed = signals[ant, start:end] * self.window
        output[ant, spec, :] = np.fft.fft(windowed)
\end{verbatim}

Loop structure prioritises clarity over maximum vectorisation. Profiling showed FFT computation dominates (0.1–0.2 ms per 256-point transform) while loop overhead remains negligible.

\subsection{Delay Engine: Geometric Corrections}

The Delay Engine computes phase corrections from antenna geometry. Setting the phase centre triggers delay calculation using vectorised matrix multiplication:

\begin{verbatim}
path_diff = ant_positions @ direction
path_diff -= path_diff[0]  # relative to antenna 0
current_delays = path_diff / c
\end{verbatim}

Applying delays uses broadcasting:

\begin{verbatim}
for ant in range(n_ants):
    phases = 2 * np.pi * channel_freqs * current_delays[ant]
    phasor = np.exp(-1j * phases)
    corrected[ant, :, :] *= phasor[np.newaxis, :]  # broadcast
\end{verbatim}

Phase correction consumes under 5\% of runtime, making further optimisation unnecessary.

\subsection{X-Engine: Correlation and Integration}

Initialisation computes baseline list (autocorrelations first, then cross-correlations) and integration parameters. Correlating a single spectrum:

\begin{verbatim}
for bl_idx, (i, j) in enumerate(baselines):
    if i == j:
        vis[bl_idx, :] = np.abs(channelized[i, :])**2
    else:
        vis[bl_idx, :] = channelized[i, :] * np.conj(channelized[j, :])
\end{verbatim}

Accumulation sums contributions until integration completes, then averages and resets. Autocorrelations are forced to real values to eliminate numerical noise from finite precision.

\subsection{Output Formatting and Configuration}

Format-specific writers handle NumPy, HDF5, and FITS outputs. HDF5 adds metadata through attributes; FITS separates complex data into real/imaginary HDUs. Missing libraries trigger graceful fallback to NumPy format.

The \texttt{CorrelatorConfig} dataclass uses type annotations and \texttt{\_\_post\_init\_\_} validation:

\begin{verbatim}
@dataclass
class CorrelatorConfig:
    n_ants: int = 4
    n_channels: int = 256
    ...
    
    def __post_init__(self):
        if not is_power_of_two(self.n_channels):
            raise ValueError("n_channels must be power of 2")
\end{verbatim}

YAML serialization converts dataclass fields to dictionaries, handling NumPy arrays by converting to nested lists.
